% ---------------------------------------------------------------------------
% Datos del grupo: se completan aquí una sola vez (carátula, fichas y conclusiones los usan).
% ---------------------------------------------------------------------------
\newcommand{\nombreRody}{Rody Sebastian Vilchez Marin}
\newcommand{\codigoRody}{U202216562}
\newcommand{\nombreDayana}{Dayana Kety Gómez Rodriguez}
\newcommand{\codigoDayana}{U202311495}
\newcommand{\nombreTercero}{Julio Cesar Meza Alfaro}
\newcommand{\codigoTercero}{U202218912}
\newcommand{\coordinador}{\nombreRody}

% ---------------------------------------------------------------------------
% Marcas de trabajo pendiente. Antes de entregar, buscar "\pendiente" y "\verificar":
% el PDF final no debe tener ninguna marca roja ni naranja.
% ---------------------------------------------------------------------------
\newcommand{\pendiente}[1]{\textcolor{BrickRed}{\itshape [Pendiente: #1]}}
\newcommand{\verificar}[1]{\textcolor{BurntOrange}{\itshape [Verificar en el PDF: #1]}}

% ---------------------------------------------------------------------------
% Tarjeta por paper: la tabla que pide el enunciado (sección e), en formato compacto.
%
%   \tarjeta{clave-bibtex}{título corto para el caption}{
%     titulo={...}, autores={...}, fuente={año · venue}, link={\url{...}},
%     motivacion={...}, problema={...}, propuesta={...}, algoritmo={\paso{1}... \paso{2}...},
%     stack={...}, aporte={...}, integrante={\nombreX},
%     opinion={...}            % la escribe el integrante; mientras falte, sale la marca roja
%   }
% ---------------------------------------------------------------------------
\pgfkeys{
  /tarjeta/.is family, /tarjeta,
  titulo/.store in=\tjTitulo,
  autores/.store in=\tjAutores,
  fuente/.store in=\tjFuente,
  link/.store in=\tjLink,
  motivacion/.store in=\tjMotivacion,
  problema/.store in=\tjProblema,
  propuesta/.store in=\tjPropuesta,
  algoritmo/.store in=\tjAlgoritmo,
  stack/.store in=\tjStack,
  aporte/.store in=\tjAporte,
  integrante/.store in=\tjIntegrante,
  opinion/.store in=\tjOpinion,
}

\newcommand{\etiqueta}[1]{\textcolor{black!70}{\bfseries #1}}
\newcommand{\separador}{\par\vspace{1.2mm}\noindent\textcolor{black!25}{\rule{\linewidth}{0.4pt}}\par\vspace{0.6mm}}
\newcommand{\paso}[1]{\par\textbf{#1.}~}

\newcommand{\tarjeta}[3]{%
  % Sin espacio antes de #3: pgfkeys solo recorta un espacio y el argumento ya trae el suyo.
  \pgfqkeys{/tarjeta}{opinion={\pendiente{la escribe el integrante: beneficios o perjuicios de la
    programación concurrente en este trabajo}},#3}%
  % No es un float: la tarjeta sigue el texto y puede partirse entre páginas, así no deja huecos.
  % \needspace evita que el título de la tabla quede solo al pie de una página.
  \par\needspace{12\baselineskip}\vspace{6pt}%
  \captionof{table}{Ficha de \textcite{#1}: #2}\label{tab:#1}
  \begin{tcolorbox}[enhanced, breakable, colback=white, colframe=black!45, boxrule=0.4pt, arc=1mm,
                    left=2.5mm, right=2.5mm, top=1.5mm, bottom=1.5mm,
                    before skip=2pt, after skip=12pt,
                    fontupper=\footnotesize\setstretch{1}\setlength{\parskip}{0pt}]
    {\small\bfseries\tjTitulo}\par\vspace{0.3mm}
    \tjAutores\par
    \tjFuente\ \textperiodcentered\ \tjLink
    \separador
    \begin{tabularx}{\linewidth}{@{}Y@{\hspace{5mm}}Y@{}}
      \etiqueta{Motivación}\par \tjMotivacion
        & \etiqueta{Problema que resuelve}\par \tjProblema \\ \addlinespace[5pt]
      \etiqueta{Propuesta de solución}\par \tjPropuesta
        & \etiqueta{Algoritmos y pseudocódigo}\par \tjAlgoritmo \\
    \end{tabularx}
    \separador
    \etiqueta{Servicios, frameworks, lenguaje y patrones.} \tjStack\par\vspace{0.8mm}
    \etiqueta{Aporte de la concurrencia.} \tjAporte
    \separador
    \etiqueta{Opinión crítica (\tjIntegrante).} \tjOpinion
  \end{tcolorbox}}

% Pseudocódigo en español.
\floatname{algorithm}{Algoritmo}
\renewcommand{\algorithmicrequire}{\textbf{Entrada:}}
\renewcommand{\algorithmicensure}{\textbf{Salida:}}
\algrenewcommand\algorithmicwhile{\textbf{mientras}}
\algrenewcommand\algorithmicdo{\textbf{hacer}}
\algrenewcommand\algorithmicfor{\textbf{para}}
\algrenewcommand\algorithmicforall{\textbf{para cada}}
\algrenewcommand\algorithmicif{\textbf{si}}
\algrenewcommand\algorithmicthen{\textbf{entonces}}
\algrenewcommand\algorithmicreturn{\textbf{devolver}}
